% Draft: Chapter 1 (Introduction/Motivation) and Chapter 3 (Requirements Analysis)
% Bachelor's Thesis in Informatics, University of Bremen — Sergi Koniashvili
% Date: 2026-11-01 — Discussion basis for supervisor meeting on 2026-11-02.
% Facts exclusively from docs/thesis/THESIS_KNOWLEDGE_BASE.md (VERIFIED).
% Unverified points are marked with [TODO].
\documentclass[11pt,a4paper]{article}
\usepackage{fontspec}
\usepackage[english]{babel}
\usepackage[a4paper,margin=2.5cm]{geometry}
\usepackage{csquotes}
\usepackage{booktabs}
\usepackage{longtable}
\usepackage{array}
\usepackage{xcolor}
\usepackage{enumitem}
\usepackage[hidelinks]{hyperref}

\newcommand{\todo}[1]{\textcolor{red!70!black}{[TODO: #1]}}
\newcommand{\question}[1]{\par\smallskip\noindent\colorbox{blue!8}{\parbox{\dimexpr\linewidth-2\fboxsep}{\small\textbf{Question for Supervisors:} #1}}\par\smallskip}

\title{Draft: Motivation and Literature-Based Requirements Analysis\\[0.3em]
\large LLM-Driven Integration of Climate and Health Data}
\author{Sergi Koniashvili\\ Bachelor's Thesis in Informatics, University of Bremen}
\date{As of: 1 November 2026 (revised) -- Discussion Basis}

\begin{document}
\maketitle

\noindent\textit{Note:} This draft presents the planned structure as well as preliminary
versions of Chapter~1 and Chapter~3. Sections marked in red are not yet fully
substantiated or remain open. The provisional citation style is Author-Year (see
question in Section~\ref{sec:questions}).

\section*{Planned Thesis Outline}
\begin{enumerate}[nosep]
  \item Introduction (Motivation, Problem Statement, Objectives, Outline)
  \item Foundations (Climate and vector-borne diseases; Data sources:
        case numbers, reanalysis climate data, population; ETL; LLMs)
  \item Requirements Analysis (Derivation from exemplary publications)
  \item Related Work and Existing Tools \todo{Verify tools from official documentation}
  \item Concept and Architecture
  \item Implementation
  \item Evaluation (Case Studies CS1--CS3)
  \item Discussion, Limitations, and Outlook
  \item Conclusion
\end{enumerate}

% ---------------------------------------------------------------------------
\section{Introduction}

\subsection{Motivation}

The interplay between climate variability and vector-borne diseases such as dengue
is increasingly investigated using data-driven models. A recent example is the
study by Suttirat et al.\ (2026), which links monthly case counts of dengue
haemorrhagic fever (DHF) across 77 provinces of Thailand over 20 years
(2003--2022) with climate, remote sensing, and socioeconomic data, analyzing them
using XGBoost and SHAP (pp.~5--8). Lowe et al.\ (2021) similarly investigate
the impact of hydrometeorological extremes and urbanisation on dengue risk across
558 microregions of Brazil (pp.~e211--e212).

Common to both studies is that actual statistical and machine-learning modeling
is preceded by an extensive, multi-step data preparation pipeline. Suttirat et al.\
describe these preprocessing steps as follows:

\begin{quote}
\enquote{Gridded climatic and remote-sensing data were spatially aggregated to
the provincial level; lagged variables and bioclimatic indices were then derived
from monthly provincial temperature and precipitation series. Socioeconomic data
were obtained directly at the provincial level: the biennially reported household
income series was linearly interpolated to fill missing years, and all
socioeconomic series were then upsampled to monthly resolution by assigning each
year's reported value (typically released at year-end) to every month of that
year.} (Suttirat et al., 2026, p.~4)
\end{quote}

These preparation procedures are not idiosyncratic to a single paper. Lowe et al.\
(2021) execute largely identical steps: they retrieve monthly notification counts
from the national surveillance system, calculate incidence per 100,000 population
using annual demographic estimates, perform area-weighted spatial aggregation of
gridded climate data (Tmin, Tmax, PDSI) to the spatial units of analysis, and
extract climate series for an extra year prior to the case observation window,
\enquote{to allow for a lag period before the first dengue observation in
January, 2001} (p.~e211). Suttirat et al.\ likewise generate time-lagged features
(lags of 1--3 months) for all climate and environmental covariates (p.~7).
Table~\ref{tab:steps} juxtaposes the processing steps of both studies.

Furthermore, while the primary datasets are frequently open-access, they are not
immediately analysis-ready. Lowe et al.\ rely entirely on open data (p.~e218), yet
these originate from five heterogeneous sources with disparate formats and temporal
resolutions (monthly, annual, static, 0.5° spatial grid); collating the database
is listed as a dedicated contributor task (p.~e218). Regarding health surveillance
data, Clarke et al.\ (2024) observe in their description of the OpenDengue database:

\begin{quote}
\enquote{Dengue case data exists in multiple formats from a wide variety of
sources that require various processing methodologies. [\ldots] Locating,
extracting, processing and standardising data all takes time but is essential to
enable reuse and re-analysis.} (Clarke et al., 2024, p.~2)
\end{quote}

OpenDengue resolves the challenge of aggregating disease surveillance data globally.
However, joining these epidemiological records with meteorological and climate
reanalysis data—explicitly cited by Clarke et al.\ as a key downstream use
case (pp.~1, 9)—remains the responsibility of individual researchers.

Finally, executing this data engineering pipeline requires substantial programming
expertise: Suttirat et al.\ implemented data preparation in Python (pandas) and
computed bioclimatic indices using the R package \texttt{dismo} (pp.~5, 8). For
domain specialists without software engineering backgrounds, such as public health
officials and clinical epidemiologists, this technical overhead poses a significant
barrier \todo{Cite literature documenting programming/data-science skill bottlenecks in public health}.
Existing software tools address these steps only in fragments
\todo{Related work: Copernicus CDS, Google Earth Engine, DHIS2 Climate \& Health, IRI ENACTS -- substantiate from official documentation}.

\subsection{Problem Statement and Objectives}

The objective of this thesis is to develop an automated software platform that
relieves researchers of the data preparation pipeline preceding model estimation—source
discovery, cleaning, harmonisation, integration, and documentation—without requiring
programming expertise, while maintaining full transparency and auditability, allowing
researchers to focus entirely on domain modeling. Users describe their required
dataset (disease, geographic region, time horizon, meteorological variables) via an
interactive form or in natural language; a Large Language Model (LLM) parses the
unstructured prompt into a structured integration request. The platform
selects appropriate authoritative sources, harmonises resolutions, joins records into
an analysis-ready tabular dataset, derives relevant epidemiological and meteorological
features (incidence, lagged variables), and audits every transformation step.

The role of the LLM is deliberately limited. It translates a natural-language
request into the fields of the form, which can be checked and changed before the
run, explains the result in plain language, and supports the chat assistant with
domain questions and form suggestions. Retrieval, cleaning, harmonisation, joining,
calculations and export, in contrast, are carried out exclusively by deterministic
program components; the same request therefore yields the same table. The LLM's
output still has to be checked: in an own test, the model described a
non-significant correlation (lag~3, $n = 9$) as the strongest association. Since
then it receives the significance flag and the number of value pairs for every
correlation and may only name significant associations as findings. How reliably
the LLM components work will be examined in the evaluation
\todo{define criteria, e.g.\ share of correctly filled form fields}.

Crucially, the system is explicitly \emph{not} a predictive modeling tool: it prepares
standardised input data for downstream analytical models. Whether it genuinely
reduces manual effort will be measured empirically in the evaluation (Chapter~7,
Case Study CS1: constructing an identical target dataset manually via Python/pandas
versus via the platform; quantifying elapsed time, manual intervention count, and
numerical concordance).

\subsection{Structure of the Thesis}
\todo{Complete after finalising the thesis outline}

% ---------------------------------------------------------------------------
\setcounter{section}{2}
\section{Requirements Analysis}

\subsection{Methodological Approach}

Functional and non-functional requirements are derived systematically from exemplary
peer-reviewed studies in climate-health modeling. For each study, the data preparation
steps executed prior to statistical modeling are extracted from the methods section.
Every step indispensable for generating an analysis-ready dataset is translated into
a concrete system requirement or explicitly scoped out with scientific justification.
The mapping of Preparation Step $\rightarrow$ Source Publication (with page reference)
$\rightarrow$ System Requirement is documented in a traceability matrix (Table~\ref{tab:trace}).

The current empirical foundation comprises three publications:
\begin{itemize}[nosep]
  \item \textbf{Suttirat et al.\ (2026)}: Dengue, Thailand, 77 provinces,
        2003--2022, XGBoost/SHAP -- primary validation case study;
  \item \textbf{Lowe et al.\ (2021)}: Dengue, Brazil, 558 microregions,
        2001--2019, Bayesian spatio-temporal modeling -- establishes the generality
        of the data preparation pipeline;
  \item \textbf{Clarke et al.\ (2024)}: OpenDengue database (Data Descriptor) --
        specifies requirements for provenance tracking, versioning, and cross-source
        comparability of epidemiological data.
\end{itemize}

\subsection{Data Preparation Steps in the Exemplary Publications}

\begin{longtable}{@{}p{0.27\linewidth}p{0.33\linewidth}p{0.33\linewidth}@{}}
\caption{Data preparation preceding model estimation (Excerpt)}\label{tab:steps}\\
\toprule
Step & Suttirat et al.\ (2026) & Lowe et al.\ (2021) \\
\midrule
\endfirsthead
\toprule
Step & Suttirat et al.\ (2026) & Lowe et al.\ (2021) \\
\midrule
\endhead
Retrieve case counts & Monthly DHF cases, Bureau of Epidemiology (p.~5) &
  Monthly notification records, SINAN/DATASUS (p.~e211) \\
Harmonise administrative units & Bueng Kan merged into Nong Khai (p.~5) & -- \\
Retrieve climate variables & TerraClimate: precipitation, Tmin, Tmax, etc.\ (p.~5) &
  CRU TS 4.04: Tmin, Tmax, PDSI (p.~e211) \\
Spatial aggregation & Gridded raster averaged to provinces, method unspecified (p.~5) &
  Area-weighted mean using \texttt{exactextractr} (p.~e211) \\
Lagged variables & Lags 1--3 months (p.~7) & Lags 0--6 months (DLNM, p.~e212) \\
Buffer for lags & -- & Climate data retrieved 1 year prior to case series (p.~e211) \\
Incidence per 100,000 & Yes (p.~7) & Yes, with annual population estimates (p.~e211) \\
Annual $\rightarrow$ Monthly & Annual value assigned to every month (pp.~4--5) &
  Not described \\
Interpolate missing years & Household income linearly interpolated (p.~5) & -- \\
Derived indices & 19 bioclimatic indicators (\texttt{dismo}, p.~5) & -- \\
Tools & Python/pandas, R (pp.~5, 8) & R, \texttt{exactextractr}, R-INLA (pp.~e211--e212) \\
\bottomrule
\end{longtable}

\subsection{Derived Requirements (Traceability Matrix)}

\begin{longtable}{@{}p{0.25\linewidth}p{0.19\linewidth}p{0.38\linewidth}p{0.14\linewidth}@{}}
\caption{Mapping: Preparation Step $\rightarrow$ Publication $\rightarrow$ Requirement}\label{tab:trace}\\
\toprule
Step / Observation & Evidence & Requirement & Status \\
\midrule
\endfirsthead
\toprule
Step / Observation & Evidence & Requirement & Status \\
\midrule
\endhead
Incidence per 100,000 as outcome & Suttirat p.~7; Lowe p.~e211 &
  FR-17: Integrate population data and compute incidence per 100,000 & Implemented \\
Annual population, monthly case counts & Suttirat pp.~4--5 (annual value assigned to every month); Lowe: method not described &
  FR-17: linearly interpolate annual values to months -- \emph{own design decision}, so that incidence does not jump at year boundaries; the difference from Suttirat is discussed in Chapter~8 & Implemented \\
Lagged climate features & Suttirat p.~7 (1--3); Lowe p.~e212 (0--6) &
  FR-23: Lagged feature columns 1--6 in integrated dataset & Implemented \\
Climate data prior to case start & Lowe p.~e211 &
  FR-23: Lag buffer query to prevent data loss in initial periods & Implemented \\
Tmin/Tmax rather than only mean & Suttirat p.~5; Lowe p.~e211 &
  FR-24: Tmin/Tmax aggregated as mean of daily extrema & Implemented \\
Disparate temporal resolutions & Suttirat pp.~4--5 &
  FR-20: Temporal resolution harmonisation \todo{Discuss difference: Suttirat repeats annual stock values, FR-20 preserves sums of count variables} & Implemented \\
Provenance must be traceable; cite versions & Clarke pp.~2, 4, 12 &
  FR-25: Export metadata package (sources, licences, version, audit steps) & Implemented \\
Reporting coverage varies across countries & Clarke pp.~1, 8 &
  FR-26: Detect temporal data gaps with diagnostic reasons, do not impute automatically & Implemented \\
Assess comparability between sources & Clarke p.~2 &
  FR-27: Empirically compare two sources over overlapping window & Implemented \\
Exploratory correlation assessment & Lags: Suttirat p.~7; Lowe p.~e212 &
  FR-19: Cross-correlations for lags 0--3 (Pearson, Spearman) & Implemented \\
\bottomrule
\end{longtable}

\noindent Regarding FR-19, an important methodological boundary must be stated:
bivariate linear correlations do not capture non-linear threshold dynamics, such as
the empirical thresholds identified by Suttirat et al.\ (approximately 21\,°C for
Tmin with lag~1 and 32\,°C for Tmax with lag~3, p.~12); Lowe et al.\ emphasize that
climate impacts on transmission are \enquote{complex, non-linear, and often delayed}
(p.~e211). Consequently, correlation coefficients serve strictly for exploratory data
verification, not as predictive model specifications.

\subsection{Requirements Without Direct Literature Origin}

The initial requirements catalogue (FR-01 to FR-22) was formulated prior to the
in-depth literature extraction. While core elements are directly corroborated by
literature (FR-17, FR-19, FR-20, see above), other requirements stem from the overarching
usability goal of enabling access \enquote{without programming skills} (natural-language
parsing FR-01, plain-language explanations FR-09, interactive Copilot assistant FR-13,
interactive world map FR-22, PDF synthesis report FR-21) and from supervisor guidelines
(e.g., FR-03 to FR-08).
\question{Should these requirements be justified primarily under Usability / Zero-Programming objectives, or streamlined from the core engineering catalogue?}

\subsection{Scoping and Delimitation}

The following data preparation steps documented in the literature are deliberately
not implemented; they are discussed as limitations in Chapter~8:
\begin{itemize}[nosep]
  \item Subnational administrative units and spatial raster averaging
        (Suttirat p.~5; Lowe p.~e211) -- the platform currently operates with a
        representative geographic coordinate per country;
  \item Administrative boundary harmonisation over time (Suttirat p.~5);
  \item Bioclimatic indicators (e.g., WorldClim bio1--bio19) and satellite remote sensing (Suttirat p.~5);
  \item Automatic cross-source gap filling: case reporting figures from different
        surveillance systems rely on divergent case definitions—OpenDengue records
        all clinically suspected and confirmed cases (Clarke p.~5), whereas Suttirat
        et al.\ focus strictly on DHF (p.~5)—such that uncoordinated merging would
        distort epidemiological comparability (A20). The platform detects gaps (FR-26)
        and provides explicit, non-destructive source comparison (FR-27).
\end{itemize}

% ---------------------------------------------------------------------------
\section*{Open Questions for the Supervisor Meeting}\label{sec:questions}
\begin{enumerate}[nosep]
  \item Does the motivational argumentation chain (Data-driven models $\rightarrow$
        Laborious, recurring data preparation $\rightarrow$ Programming barrier
        $\rightarrow$ Automated platform) align with your expectations?
  \item Are two empirical modeling papers and one surveillance data descriptor sufficient,
        or should additional diseases (e.g., malaria) be incorporated?
  \item Is the traceability matrix considered an appropriate format for requirements derivation in informatics?
  \item Proposed research question: \enquote{To what extent can an LLM-supported
        system automate the integration of heterogeneous climate and health data
        while making the correctness, reproducibility and traceability of the data
        preparation verifiable?} Does this direction fit?
  \item Which citation style is preferred (Author-Year vs.\ numeric)?
  \item Thesis structure: Is the placement of \enquote{Related Work} after the Requirements Analysis appropriate?
\end{enumerate}

\section*{References (Draft)}
\begin{small}
\begin{description}[nosep,leftmargin=1em,labelindent=0pt]
\item Clarke, J., Lim, A., Gupte, P., Pigott, D.\,M., van Panhuis, W.\,G., \&
  Brady, O.\,J. (2024). A global dataset of publicly available dengue case count
  data. \emph{Scientific Data}, 11, 296.
  \url{https://doi.org/10.1038/s41597-024-03120-7}
\item Lowe, R., Lee, S.\,A., O'Reilly, K.\,M., Brady, O.\,J., Bastos, L.,
  Carrasco-Escobar, G., et al.\ (2021). Combined effects of hydrometeorological
  hazards and urbanisation on dengue risk in Brazil: a spatiotemporal modelling
  study. \emph{The Lancet Planetary Health}, 5(4), e209--e219.
  \url{https://doi.org/10.1016/S2542-5196(20)30292-8}
\item Suttirat, P., Chadsuthi, S., Aekthong, S., Rocklöv, J., Bicout, D.\,J.,
  Haddawy, P., Su Yin, M., Lawpoolsri, S., \& Modchang, C. (2026). Machine
  learning reveals temperature as a key predictor of dengue risk across
  Thailand's provinces: A 20-year analysis. \emph{PLoS Neglected Tropical
  Diseases}, 20(7), e0014590.
  \url{https://doi.org/10.1371/journal.pntd.0014590}
\end{description}
\end{small}

\end{document}
